\section*{Capítulo \ref{ch:g12:polymers} --- Polímeros}

\begin{solution}{exo:g12:polymers:1}
El propeno, \ce{CH2=CH-CH3}.
\end{solution}

\begin{solution}{exo:g12:polymers:2}
\ce{-CF2-CF2-}; un \omterm{def:g12:polymers:addition-polymer}{polímero de adición}: el \omterm{def:g12:polymers:polymer}{monómero} tiene un enlace
\ce{C=C} y no se forma ningún otro producto.
\end{solution}

\begin{solution}{exo:g12:polymers:3}
De adición: polietileno, PVC, poliestireno. De condensación: PET (un
poliéster), nailon (una poliamida).
\end{solution}

\begin{solution}{exo:g12:polymers:4}
La celulosa, el almidón, el caucho del látex y las proteínas. El nailon y
el PVC son sintéticos.
\end{solution}

\begin{solution}{exo:g12:polymers:5}
Cada enlace usa un grupo de cada \omterm{def:g12:polymers:polymer}{monómero}; con dos grupos, un \omterm{def:g12:polymers:polymer}{monómero}
puede enlazarse por ambos lados y la cadena sigue creciendo. Un \omterm{def:g12:polymers:polymer}{monómero}
con un solo grupo detendría la cadena.
\end{solution}

\begin{solution}{exo:g12:polymers:6}
$M(\ce{C2H3Cl}) = 24.0 + 3.0 + 35.5 = \qty{62.5}{g/mol}$;
$n = 125000/62.5 = 2000$.
\end{solution}

\begin{solution}{exo:g12:polymers:7}
$M(\ce{C8H8}) = \qty{104.0}{g/mol}$; $M = 2500 \times 104.0 =
\qty{2.60e5}{g/mol}$.
\end{solution}

\begin{solution}{exo:g12:polymers:8}
El grupo ácido de una \omterm{def:g7:atoms-and-molecules:molecule}{molécula} forma un \omterm{def:g11:functional-groups:carboxylic-acid}{éster} con el grupo \omterm{def:g11:functional-groups:alcohol}{alcohol} de la
siguiente, que todavía tiene un grupo ácido libre: la cadena crece.
\omterm{def:g12:polymers:repeat-unit}{Unidad de repetición} \ce{-O-CH(CH3)-CO-}; en cada enlace se libera agua.
\end{solution}

\begin{solution}{exo:g12:polymers:9}
Las películas de las bolsas están hechas de cadenas ramificadas, que no
pueden empaquetarse: el material es amorfo, blando y transparente. Las
botellas de leche están hechas de cadenas casi lineales, que se
empaquetan en regiones cristalinas: más densas, más rígidas, y las
regiones cristalinas dispersan la luz.
\end{solution}

\begin{solution}{exo:g12:polymers:10}
Sus cadenas están unidas por muchos entrecruzamientos covalentes. Al
calentar no se consigue que las cadenas se deslicen; se rompen enlaces y
el material se destruye en lugar de fundirse.
\end{solution}

\begin{solution}{exo:g12:polymers:11}
$M(\ce{C12H22N2O2}) = 144.0 + 22.0 + 28.0 + 32.0 = \qty{226.0}{g/mol}$;
$n = 22600/226.0 = 100$.
\end{solution}

\begin{solution}{exo:g12:polymers:12}
\[
  \ce{H2N(CH2)6NH2 + HOOC(CH2)4COOH -> H2N(CH2)6NHCO(CH2)4COOH + H2O} .
\]
El producto todavía lleva un grupo \omterm{def:g11:functional-groups:amine}{amina} libre en un extremo y un grupo
ácido libre en el otro: cada uno puede reaccionar con otro \omterm{def:g12:polymers:polymer}{monómero}.
\end{solution}

\begin{solution}{exo:g12:polymers:13}
Dos \omterm{def:g7:atoms-and-molecules:molecule}{moléculas} de agua por \omterm{def:g12:polymers:repeat-unit}{unidad de repetición}. $n = 1000/226.0 =
\qty{4.42}{mol}$ de \omterm{def:g12:polymers:repeat-unit}{unidades de repetición}, así que \qty{8.85}{mol} de
agua:
$8.85 \times 18.0 = \qty{159}{g}$.
\end{solution}

\begin{solution}{exo:g12:polymers:14}
El caucho natural caliente está formado por cadenas que se deslizan unas
sobre otras. Los puentes de azufre unen las cadenas: se estiran, pero
recuperan su forma, y ya no fluyen. Al estar entrecruzado, un neumático
no puede fundirse y volver a moldearse.
\end{solution}

\begin{solution}{exo:g12:polymers:15}
$M(\ce{C6H10O5}) = \qty{162.0}{g/mol}$; $n = \num{1.6e6}/162.0 = 9900$
unidades; longitud $9900 \times 0.5 = \qty{4900}{nm}$, unos
\qty{5}{\micro m}.
\end{solution}

\begin{solution}{pb:g12:polymers:1}
\textbf{1.} El etano-1,2-diol (dos grupos \omterm{def:g11:functional-groups:alcohol}{alcohol}) y el ácido
benceno-1,4-dicarboxílico (dos grupos \omterm{def:g11:functional-groups:carboxylic-acid}{ácido carboxílico}).

\textbf{2.} Un enlace \omterm{def:g11:functional-groups:carboxylic-acid}{éster}; se libera agua.

\textbf{3.} Un \omterm{def:g12:polymers:addition-polymer}{polímero de condensación}.

\textbf{4.} $M(\ce{C2H6O2}) = \qty{62.0}{g/mol}$; $M(\ce{C8H6O4}) =
\qty{166.0}{g/mol}$.

\textbf{5.} $62.0 + 166.0 - 2 \times 18.0 = \qty{192.0}{g/mol}$; y
$10 \times 12.0 + 8 \times 1.0 + 4 \times 16.0 = 192.0$.

\textbf{6.} $n = 38400/192.0 = 200$.

\textbf{7.} Dos enlaces \omterm{def:g11:functional-groups:carboxylic-acid}{éster} por \omterm{def:g12:polymers:repeat-unit}{unidad de repetición}: unos 400.

\textbf{8.} Sus cadenas son mayoritariamente lineales, y partes de ellas
se alinean en regiones ordenadas.

\textbf{9.} Es un \omterm{def:g8:plastics:thermoplastic}{termoplástico}: sus cadenas no están entrecruzadas y se
deslizan unas sobre otras en caliente.

\textbf{10.} $300/0.80 = \qty{375}{g}$.

\textbf{11.} $375/25 = 15$ botellas.

\textbf{12.} Los tapones, las etiquetas, la suciedad y los recortes se
eliminan o se pierden durante la clasificación, el lavado y el hilado.

\textbf{13.} Prevenir los residuos (principio 1): un objeto usado se
convierte en \omterm{def:g5:raw-materials-and-recycling:raw-material}{materia prima}.

\textbf{14.} \ce{C10H8O4 + 2H2O -> C2H6O2 + C8H6O4}.

\textbf{15.} $1000/192.0 = \qty{5.21}{mol}$.

\textbf{16.} Diol: $5.21 \times 62.0 = \qty{323}{g}$; ácido:
$5.21 \times 166.0 = \qty{865}{g}$.

\textbf{17.} $2 \times 5.21 \times 18.0 = \qty{188}{g}$ de agua;
$1000 + 188 = 323 + 865 = \qty{1188}{g}$: la masa se conserva.

\textbf{18.} Los \omterm{def:g12:polymers:polymer}{monómeros} pueden purificarse y dar PET nuevo tan bueno
como el original, mientras que la fusión acorta un poco las cadenas cada
vez.

\textbf{19.} Un forro polar necesita 15 botellas de \qty{25}{g}.
\end{solution}
